\documentclass[openany,oneside,a4paper,10pt]{extarticle}
% acesso a links da internet
%\usepackage[pdftext]{hyperref}
\usepackage{hyperref}
\usepackage{setspace}
\usepackage[lmargin=2.5cm,tmargin=2.5cm,rmargin=2.5cm,bmargin=2.5cm]{geometry}
%\usepackage[margin=2cm,headheight=1.5cm, includeheadfoot]{geometry}
%\usepackage[parfill]{parskip}
% DEFININDO OS PADRÕES DO EVENTO DA UENF
% VII Colóquio Internacional de Cognição e Linguagem

\usepackage{lipsum} % Textos aleatórios...

\usepackage{float}

\usepackage{graphicx}
\usepackage{subfig}
\usepackage{longtable}  % <-- Novo! Para tabelas com mais de uma página
\usepackage[font=scriptsize,labelfont=bf,justification=centering]{caption}

\usepackage[brazil]{babel}
\usepackage[T1]{fontenc}% Selecao de codigos de fonte.
\usepackage{helvet}
\renewcommand{\familydefault}{\sfdefault}

\usepackage[utf8]{inputenc}
\usepackage{tabularx}   % Para a tabela se ajustar à largura do texto

\usepackage{booktabs}   % Para linhas de tabela (toprule, midrule, bottomrule)
%\usepackage{array}      % Para >{\raggedright\arraybackslash}

\usepackage[table]{xcolor} % Importa o pacote para cores na tabela

\usepackage{array} % Pacote necessário para alinhamento vertical

% No pre{\^a}mbulo do seu documento, certifique-se de ter os pacotes:
\usepackage{amsmath}
\usepackage{amssymb}
\usepackage{amsthm}
% E tamb{\'e}m defina os ambientes de teorema, se ainda n{\~a}o o fez:
\newtheorem{theorem}{Teorema}
\newtheorem{definition}{Defini{\c c}{\~a}o}

% Cabeçalho e rodapé
\usepackage{tikz}
\usepackage{fancyhdr}
\pagestyle{fancy} % Define 'fancy' como o estilo padrão para todas as páginas

% --- Configurações GLOBAIS de Estilo ---
\renewcommand{\headrulewidth}{0pt}     % Remove a linha do cabeçalho
\renewcommand{\footrulewidth}{0.4pt}    % ADICIONA A LINHA ACIMA DO RODAPÉ

% --- Conteúdo para as páginas 2 em diante (estilo PADRÃO) ---
\fancyhf{} % Limpa os campos de cabeçalho e rodapé antes de redefinir
\fancyfoot[L]{\footnotesize ISSN: 16799844 – InterSciencePlace – International Scientific Journal}
\fancyfoot[R]{\footnotesize Page \thepage}

% --- Configuração ESPECIAL apenas para a PRIMEIRA página ---
\fancypagestyle{meuEstilo}{%
  \fancyhf{} % Limpa tudo para este estilo específico
  % Note: As larguras das linhas (\headrulewidth e \footrulewidth) são herdadas da configuração global
  %\fancyhead[C]{
  %  \centering
  %  \includegraphics[width=21cm]{cabecalhoArtigoUNEF.png}
  %}
  \fancyhead{}
  \renewcommand{\footrulewidth}{0pt}    % ADICIONA A LINHA ACIMA DO RODAPÉ
  %\fancyfoot[C]{\footnotesize Av. Alberto Lamego, 2000 - Parque Califórnia - Campos dos Goytacazes/ RJ\\
  %Tel.: (22) 2739.7186 - Fax: (22) 2739.7281 - email: pgclcch@uenf.br}
}
% Pacotes de citações BibLaTeX
% ----------------------------------------------------------
\usepackage[style=abnt,
	backend=biber,
	citecounter=true,
	backrefstyle=three,
	url=true,
	maxbibnames=99,
    mincitenames=1,
    maxcitenames=2,
    backref=false,
    hyperref=true,
    giveninits=true,
    uniquename=false,
    uniquelist=false]{biblatex}

% Colocar et al em itálico
\usepackage{xpatch}
\xpatchbibmacro{name:andothers}{%
  \bibstring{andothers}%
}{%
  \bibstring[\emph]{andothers}%
}{}{}

% Arquivo de entrada das referências bibliográficas
\addbibresource{referencias.bib}

% Norma NBR10.520/2023 - Sobrenomes em minúsculas
\renewcommand*{\mkbibnamefamily}[1]{#1}%
 % Norma NBR10.520/2023 - Sobrenomes em minúsculas
%\renewcommand*{\mkbibnamefamily}[1]{#1}%

%fullcite com todos autores e não como cite
\makeatletter
\newcommand{\tempmaxup}[1]{\def\blx@maxcitenames{99}#1}
\makeatother

\DeclareCiteCommand{\fullcite}[\tempmaxup]
{\usebibmacro{prenote}}
{\usedriver
	{}
	{\thefield{entrytype}}}
{\multicitedelim}
{\usebibmacro{postnote}}

\setlength\bibitemsep{12pt}

%Centralizando o Caption das imagens e tabelas
%\usepackage[justification=centering]{caption}
\usepackage[table]{xcolor}

\title{\input{Título-autores-resumo-palavras/titulo}}

\author{
        \input{ Título-autores-resumo-palavras/autores}
        }
\usepackage{quoting}
\usepackage{abstract}
\usepackage{indentfirst}
\setlength{\parindent}{30pt} % padrão 15pt.
%\setlength{\parskip}{1cm plus 4mm minus 3mm}

%\usepackage{sectsty}
%\sectionfont{\fontsize{12}{15}\selectfont}
%\subsectionfont{\fontsize{12}{15}\selectfont}

\usepackage{outlines}

% Definição da Macro para Citação Longa (mais de 3 linhas)
\newcommand{\citelonga}[1]{
    \begin{quoting}[rightmargin=0cm,leftmargin=4cm]
        \small
        \begin{singlespace}
            \noindent #1
        \end{singlespace}
    \end{quoting}
}

\usepackage{caption}
\usepackage{tabularx}
\usepackage{ltablex}
\usepackage{newfloat} % Pacote para criar novos tipos de elementos flutuantes

% --- Criando o elemento Independente "Quadro" ---
\DeclareFloatingEnvironment[
    fileext=loq,
    listname={Lista de Quadros},
    name=Quadro,
    placement=htbp,
]{quadroenv} 

% --- Configuração das Legendas (Padronização ABNT) ---
\captionsetup[quadroenv]{
    labelsep=endash,
    justification=centering,
    font=small,
    labelfont=bf,
    skip=10pt
}
% Força a numeração independente para o novo ambiente
\makeatletter
\@addtoreset{quadroenv}{section} % Opcional: reinicia por seção, se desejar
\makeatother
\captionsetup[table]{
    labelsep=endash,
    justification=centering,
    font=small,
    labelfont=bf,
    skip=10pt
}

% Defina aqui o número da página inicial fornecido pela revista
\input{pag_inical}

\begin{document}

%\begin{titlepage}
\begingroup
%    % --- Novo Bloco com TikZ (A SOLUÇÃO) ---
%\begin{tikzpicture}[remember picture, overlay]
%    \node[anchor=north, yshift=0cm] at (current page.north) {
%        \includegraphics[width=\paperwidth,height=3.8cm,keepaspectratio]{cabecalhoArtigoUENF.png}
%    };
%\end{tikzpicture}

% --- Novo Bloco de Cabeçalho Editável ---
\begin{tikzpicture}[remember picture, overlay]
    % Insere a imagem de fundo (agora sem o texto estático)
    \node[anchor=north, yshift=0cm] at (current page.north) {
        \includegraphics[width=\paperwidth,height=3.8cm,keepaspectratio]{cabecalhoArtigoUENF.png}
    };

    % Insere o texto editável sobreposto
    % O posicionamento (x,y) pode precisar de pequenos ajustes dependendo da sua imagem
    \node[anchor=west, text width=12cm, font=\sffamily\normalfont, color=blue!40!black] at ([shift={(8.2cm,-2.1cm)}]current page.north west) { \input{edicao}};
\end{tikzpicture}

    % 2. Garante que esta página não tem cabeçalho/rodapé padrão
    \thispagestyle{meuEstilo}

    % 3. Adiciona um espaço vertical para empurrar o título para baixo
    \vspace*{1cm} % <-- AJUSTE ESTE VALOR PARA POSICIONAR O TÍTULO

    % 4. Insere o TÍTULO manualmente
    \begin{center}
        \LARGE\bfseries
        \input{Título-autores-resumo-palavras/titulo}
    \end{center}
    \vspace{1cm} % Espaço entre título e autores

    % 5. Insere os AUTORES manualmente
    % O \noindent evita a indentação do primeiro autor
    \noindent\input{Título-autores-resumo-palavras/autores}
    \vspace{1.5cm} % Espaço entre autores e resumo

    % 6. Insere o RESUMO e PALAVRAS-CHAVE manualmente
    \setlength{\parindent}{0pt} % Garante que não há indentação no resumo
    \noindent\textbf{\textit{Resumo - }}
    \input{Título-autores-resumo-palavras/resumo}
    \par\vspace{1\baselineskip}
    \noindent\textbf{Palavras-chave: } \input{Título-autores-resumo-palavras/palavras-chave}
    \par\vspace{1\baselineskip}
    \noindent\textbf{\textit{Abstract - }}
    \input{Título-autores-resumo-palavras/abstract}
    \par\vspace{1\baselineskip}
    \noindent\textbf{Keywords: } \input{Título-autores-resumo-palavras/keywords}

\endgroup
%\end{titlepage}
% --- FIM DO NOVO BLOCO ---

%**************
% Introdução **
%**************
\section{Introdução}
\onehalfspacing
\input{secoes/introducao}


%*************************
% Fundamentação Teórica **
%*************************
\section{Fundamentação Teórica}
\onehalfspacing
\input{secoes/fund_teorica}

%***************
% Metodologia **
%***************
\section{Metodologia}
\onehalfspacing
\input{secoes/metodologia}

%**************
% Resultados **
%**************
\section{Resultados e Discussão}
\onehalfspacing
\input{secoes/resultados}

%************************
% Considerações Finais **
%************************
\section{Conclusão}
\onehalfspacing
\input{secoes/consideracoes_finais} 

%*** Bibliografia ***
\printbibliography

%*************
% Apêndices **
%*************
%\newpage
%\section*{Apêndice}
%\onehalfspacing
%\input{secoes/apendices} 

\end{document}
